\documentclass{../nndlcourse}

\begin{document}

\title{大项目要求}
\author{数据科学与大数据技术专业}
\date{ }
\maketitle


% ============================================================
\section{总体说明}
% ============================================================

本课程大项目采用开放式形式, 不限定为单一类型的作业. 你可以结合自己的兴趣、已有基础和团队分工, 选择网站开发、手机应用开发、论文复现与改进, 或其他与机器学习相关的实验项目. 只要项目内容充实、技术路线清楚、工作量充分且能够体现你个人或团队的真实投入, 均可作为课程大项目提交.

本项目支持\textbf{多次提交}. 你可以在学期内持续迭代自己的开源仓库, 不断补充功能、完善文档、修复问题并更新实验结果. 最终成绩将结合项目完成质量、仓库贡献情况以及答辩表现综合给出.


% ============================================================
\section{具体要求}
% ============================================================

课程大作业的具体要求如下:

\begin{enumerate}
	\item \textbf{提交截止日期}
	\begin{enumerate}
		\item 大项目的最终提交截止日期为\textbf{期末考试前一周}.
		\item 由于本项目支持多次提交, 建议你不要等到最后一次性集中完成, 而应尽早建立仓库并持续更新, 以便完整展示项目的开发过程.
	\end{enumerate}

	\item \textbf{仓库开源要求}
	\begin{enumerate}
		\item 大作业必须以\textbf{开源仓库}的形式提交, 推荐使用 GitHub 或 Gitee.
		\item 最终提交给教师的内容, 应当是你的开源仓库链接, 而不是单独打包的压缩文件.
		\item 开源仓库中应当包含项目代码、必要的说明文档、运行方式、实验结果展示, 以及你在完成项目过程中与 AI 工具交流的记录文件.
		\item 与 AI 的交流记录请以 Markdown 文件形式保存在仓库中, 便于查看你在项目实施过程中的提问、分析、尝试和修改过程.
	\end{enumerate}

	\item \textbf{团队协作规则}
	\begin{enumerate}
		\item 本项目最多支持 \textbf{3 人一组}合作完成.
		\item 如果采用组队形式, 要求每一位成员都必须对开源仓库有真实、可见的贡献, 不能出现只有个别人完成主要工作、其他成员几乎没有参与的情况.
		\item 我需要能够从仓库中清晰地看到每位贡献者的提交历史, 即 Commit History 应当能够反映各成员实际承担的工作内容.
		\item 因此, 请不要把所有工作都集中到同一个账号下提交, 也不要在最后阶段一次性代替其他成员统一上传全部内容.
	\end{enumerate}

	\item \textbf{项目选题模式}
	\begin{enumerate}
		\item 开发一个有趣的网站.
		\item 开发一个有趣的手机应用.
		\item 对经典论文进行复现, 或在复现基础上进一步改进.
		\item 任何与机器学习相关、且具有一定技术内容与完成度的实验项目.
	\end{enumerate}

	\item \textbf{考核与答辩}
	\begin{enumerate}
		\item 如果希望获得较高分数, 请主动申请在\textbf{期末考试前一周}进行项目答辩.
		\item 答辩时间约为 10 分钟, 你需要介绍项目的主要内容、技术路线、完成情况以及个人负责的具体模块.
		\item 若为多人合作项目, 则每位成员都需要上台讲解自己负责的工作内容, 同时也应能够说明项目整体思路与系统结构.
		\item 在完成大作业过程中, 你可以使用高级 AI 工具.
		\item 但在答辩过程中, 我会围绕项目实现细节、技术方案、关键设计与实际贡献进行提问, 并根据你的回答情况综合评定最终成绩.
	\end{enumerate}

	\item \textbf{关于高质量项目的额外说明}
	\begin{enumerate}
		\item 对于完成质量非常高、工作量非常充分且具有明显创造性的大项目, 若整体表现特别优秀, 可以仅依据大项目成绩给出平时分, 平时作业可免做.
		\item 但是, 如果你希望申请以大项目替代平时作业, 则该项目\textbf{必须参加项目答辩}.
		\item 换言之, 只有项目本身做得足够扎实, 且你能够在答辩中清楚说明自己的工作内容与技术细节, 才可能获得这一特殊安排.
	\end{enumerate}
\end{enumerate}


\end{document}
